\subsection{理论的比较}

一个理论 \(\mathcal{C}'\) 被认为强于理论 \(\mathcal{C}\) ，如果 \(\mathcal{C}\) 的所有符号都是 \(\mathcal{C}'\) 的符号， \(\mathcal{C}\) 的所有显式公理都是 \(\mathcal{C}'\) 中的定理，以及 \(\mathcal{C}\) 的方案是 \(\mathcal{C}'\) 的方案。

C4. 如果理论 \(\mathfrak{G}'\) 强于理论 \(\mathfrak{G}\) ，那么 \(\mathfrak{G}\) 的所有定理都是 \(\mathfrak{G}'\) .

设 \(R_1, R_2, \ldots, R_n\) 是 \(\mathfrak{G}\) 中的一个证明。我们将逐步证明，每一个 \(R_i\) 是 \(\mathfrak{G}'\) 中的定理 。假设这一点对先于 \(R_k\) 的关系成立。如果 \(R_k\) 是 \(\mathfrak{G}\) 的公理，则它是 \(\mathfrak{G}'\) 中的一个定理，根据假设。如果 \(R_k\) 前面有关系 \(R_i\) 和 \(R_i \Longrightarrow R_k\) ，我们已经知道 \(R_i\) 和 \(R_i \Longrightarrow R_k\) 是 \(\mathfrak{G}'\) 中的定理，从而 \(R_k\) 是 \(\mathfrak{G}'\) 中的定理 ，根据C1。因此，在每种情况下， \(R_k\) 是 \(\mathfrak{G}'\) 中的定理 ，证明完成。

如果两个理论 \(\mathfrak{G}\) 和 \(\mathfrak{G}'\) 中的每一个都比另一个更强， \(\mathfrak{G}\) 和 \(\mathfrak{G}'\) 被称为等价的。那么 \(\mathfrak{G}\) 的每一个定理都是 \(\mathfrak{G}'\) 中的一个定理，反之亦然。

C5. 设 \(\mathfrak{G}\) 设为一个理论，令 \(A_{1}, A_{2}, \ldots, A_{n}\) 为其显式公理， \(a_{1}, a_{2}, \ldots, a_{h}\) 其常量，并令 \(T_{1}, T_{2}, \ldots, T_{h}\) 为 \(\mathfrak{G}\) 中的项。假设

\[
\left(\boldsymbol {T} _ {1} \mid \boldsymbol {a} _ {1}\right) \left(\boldsymbol {T} _ {2} \mid \boldsymbol {a} _ {2}\right) \dots \left(\boldsymbol {T} _ {h} \mid \boldsymbol {a} _ {h}\right) \boldsymbol {A} _ {i} \quad (\text { 对 } i = 1, 2, \dots , n)
\]

是某个理论 \(\mathfrak{G}'\) 中的定理，并且 \(\mathfrak{G}\) 的所有符号都是 \(\mathfrak{G}'\) 的符号，并且 \(\mathfrak{G}\) 的方案是 \(\mathfrak{G}'\) 的方案。那么，若 \(A\) 是 \(\mathfrak{G}\) 中的定理 ,

\[
\left(\boldsymbol {T} _ {1} \mid \boldsymbol {a} _ {1}\right) \dots \left(\boldsymbol {T} _ {h} \mid \boldsymbol {a} _ {h}\right) \boldsymbol {A}
\]

是 \(\mathfrak{C}'\) 中的定理 .

因为 \(\mathfrak{c}'\) 比理论 \((T_1|a_1) \ldots (T_h|a_h)\mathfrak{c}\) 更强，所以可以应用C2和C4。

当我们使用这一程序来从 \(\mathfrak{C}\) 的一个定理推导 \(\mathfrak{C}'\) 中的一个定理时，我们说我们在 \(\mathfrak{C}'\) 中应用 \(\mathfrak{C}\) 的结果。直观上，公理 \(\mathfrak{C}\) 表达了 \(a_1, a_2, \ldots, a_h\) 的一些性质，而 \(A\) 表达了由这些公理推出的一个性质。如果这些对象 \(T_1, T_2, \ldots, T_h\) 在 \(\mathfrak{C}'\) 中具有由公理 \(\mathfrak{C}\) 所表达的性质，那么它们也具有性质 \(A\) .

*例如，在群论 \(\mathfrak{C}\) 中，显式公理包含两个常量 \(G\) 和 \(\mu\) （群和合成法则）。在集合论 \(\mathfrak{C}'\) 中，我们定义两个项：实数直线和实数加法。如果我们将这些项分别代入 \(G\) 和 \(\mu\) ，也就是在 \(\mathfrak{C}\) 的显式公理中作此代入，我们就在 \(\mathfrak{C}'\) 中得到定理。此外， \(\mathfrak{C}\) 和 \(\mathfrak{C}'\) 的方案和符号是相同的。因此，我们可以“将群论的结果应用于实数的加法群”。我们说，我们在集合论中为群论构造了一个模型。（注意，由于群论比集合论更强，我们也可以将集合论的结果应用于群论。）*

备注。在C5的假设下，如果理论 \(\mathfrak{C}\) 结果是矛盾的，那么同样的情况也会适用于 \(\mathfrak{C}'\) 。因为如果 \(A\) 与“非 \(A\)”都是 \(\mathfrak{C}\) 中的定理，那么 \((T_1|a_1)\ldots (T_h|a_h)A\) 和 \(\mathrm{not}(T_1|a_1)\ldots (T_h|a_h)A\) 是 \(\mathfrak{C}'\) 中的定理。* 例如，如果群论是矛盾的，那么集合论也将是矛盾的。*

